% ============================================================
% Cheat sheet, argomento: il pattern State (dispensa, capitoli 9 e 15)
% ============================================================
\argomento{State}{Dispensa: capitoli 9 e 15. Pattern comportamentale della GoF}[Pattern State]

\scheda{Carta d'identità}{
  \pintento{Permettere a un oggetto di \textbf{cambiare comportamento quando cambia il suo stato interno}: sembra che
    cambi classe. È il modo a oggetti di realizzare un \textbf{diagramma di stato}.}
  \ptesto{«si comporta in modo diverso a seconda dello stato (della fase)», «lo stesso comando ha effetti diversi», un
    \textbf{diagramma di stato} da implementare, «stati e transizioni»; gli esempi tipici: il \textbf{distributore
    automatico}, un ordine (in preparazione, spedito, consegnato), una connessione (aperta, chiusa), un'interfaccia
    grafica con più schermate.}
  \pcodice{il \textbf{contesto} ha un campo di un tipo astratto (\code{State stato}) e \textbf{delega} ogni evento allo
    stato corrente passandogli se stesso: \code{stato.inserisciMoneta(this)}; una \textbf{sottoclasse per stato}, con un
    metodo per evento; le \textbf{transizioni le fanno gli stati stessi}, chiamando \code{contesto.setState(\ldots)};
    spesso gli stati sono oggetti condivisi, creati una volta in campi statici.}
  \pdiagramma{\code{Context} con una composizione (o aggregazione) \code{1} verso la classe astratta \code{State}, e le
    sottoclassi \code{ConcreteStateA}, \code{ConcreteStateB}\ldots; la nota \code{state.handle()} sul contesto.}
  \pvarianti{\code{State} classe astratta con i metodi vuoti (un evento non previsto viene ignorato) o interfaccia;
    stati condivisi e senza campi (i dati stanno nel contesto) oppure uno stato per contesto, se lo stato deve
    ricordare qualcosa; senza pattern, gli \code{switch} annidati (\vv{statoswitch}).}
  \pconfondere{lo \textbf{Strategy} (\vv{strategy-struttura}): \textbf{stessa struttura}, ma l'oggetto lo sceglie il
    \textbf{client} e le strategie non si conoscono fra loro; nello State sono gli \textbf{stati} a decidere il
    prossimo stato, e ognuno sa in quali altri si può passare.}
}

\begin{struttura}{state-struttura}{La struttura e la corrispondenza con il diagramma di stato}
\begin{tikzpicture}
  \umlclass[2.2cm]{c}{(0,0)}{Context}{}{+ request()}
  \umlclass[2.2cm]{s}{(4.0,0)}{\textit{State}}{}{+ handle(c)}
  \umlclass[2.6cm]{a}{(2.4,-2.4)}{ConcreteStateA}{}{+ handle(c)}
  \umlclass[2.6cm]{b}{(5.6,-2.4)}{ConcreteStateB}{}{+ handle(c)}
  \draw[compos] (c.east) -- (s.west);
  \estO{(s.west)}{1}{state}
  \draw[gen] (a.north) |- ++(0,0.45) -| (s.south);
  \draw[gen] (b.north) |- ++(0,0.45) -| (s.south);
  \node[notauml, anchor=north] (n) at (0,-1.3) {state.handle(this)};
  \draw[tratt] (c.south) -- (n.north);
\end{tikzpicture}
\tcblower
\textbf{Context}: l'oggetto che il client usa; tiene lo stato corrente e i dati. \textbf{State}: un metodo per ogni
evento, vuoti. \textbf{ConcreteState}: uno per stato del diagramma; ridefinisce solo gli eventi che lì hanno un
effetto.

\renewcommand{\arraystretch}{1.15}\footnotesize
\begin{tabular}{@{}l l@{}}
  \textbf{Diagramma di stato} & \textbf{Codice} \\ \hline
  stato & sottoclasse di \code{State} \\
  evento & metodo di \code{State} \\
  guardia & \code{if} nel metodo \\
  azione & corpo del metodo \\
  transizione & \code{setState(\ldots)} \\
\end{tabular}
\end{struttura}

\begin{programma}{state-esempio}{L'esempio del corso: il distributore automatico (ridotto)}
\begin{lstlisting}[style=javacs, name=Cliente.java]
public class Cliente {
    public static void main(String[] args) {
        Distributore d = new Distributore();
        d.selezionaProdotto();
        d.inserisciMoneta();
        d.inserisciMoneta();
        d.selezionaProdotto();
        d.inserisciMoneta();
        d.prelevaProdotto();
    }
}
class Distributore {                     // Context
    static final State ATTESA = new Attesa(),
            SCELTA = new AttesaScelta(),
            EROGAZIONE = new Erogazione();
    private State stato = ATTESA;
    private int credito = 0;

    void setState(State s) { stato = s; }
    int getCredito() { return credito; }
    void setCredito(int c) { credito = c; }
    // ogni evento è delegato allo stato corrente
    void inserisciMoneta() { stato.moneta(this); }
    void selezionaProdotto() { stato.scelta(this); }
    void prelevaProdotto() { stato.prelievo(this); }
}
abstract class State {            // eventi: vuoti
    void moneta(Distributore d) { }
    void scelta(Distributore d) { }
    void prelievo(Distributore d) { }
}
class Attesa extends State {
    void moneta(Distributore d) {
        d.setCredito(d.getCredito() + 1);
        System.out.println("Credito: "
                + d.getCredito());
        if (d.getCredito() >= 2)       // guardia
            d.setState(Distributore.SCELTA);
    }
    void scelta(Distributore d) {
        System.out.println("Credito insufficiente");
    }
}
class AttesaScelta extends State {
    void moneta(Distributore d) {
        d.setCredito(d.getCredito() + 1);
        System.out.println("Credito: "
                + d.getCredito());
    }
    void scelta(Distributore d) {
        System.out.println("Erogazione...");
        d.setState(Distributore.EROGAZIONE);
    }
}
class Erogazione extends State {
    void moneta(Distributore d) {
        System.out.println("Attendi l'erogazione");
    }
    void prelievo(Distributore d) {
        d.setCredito(d.getCredito() - 2);
        System.out.println("Ecco il prodotto, "
                + "credito: " + d.getCredito());
        if (d.getCredito() >= 2)
            d.setState(Distributore.SCELTA);
        else
            d.setState(Distributore.ATTESA);
    }
}
\end{lstlisting}
\tcblower
\begin{tikzpicture}
  \umlclass[3.6cm]{d}{(0,0)}{Distributore}{\underline{\textasciitilde\ ATTESA: State}\\\underline{\textasciitilde\ SCELTA: State}\\\underline{\textasciitilde\ EROGAZIONE: State}\\- credito: int = 0}{\textasciitilde\ inserisciMoneta()\\\textasciitilde\ selezionaProdotto()\\\textasciitilde\ prelevaProdotto()\\\textasciitilde\ setState(s: State)}
  \umlclass[2.6cm]{s}{(4.2,0)}{\textit{State}}{}{\textasciitilde\ moneta(d)\\\textasciitilde\ scelta(d)\\\textasciitilde\ prelievo(d)}
  \node[classe1=1.4cm] (a) at (2.2,-3.0) {Attesa};
  \node[classe1=1.9cm] (b) at (3.95,-3.0) {AttesaScelta};
  \node[classe1=1.8cm] (e) at (5.8,-3.0) {Erogazione};
  \draw[compos] (d.east |- s.west) -- (s.west);
  \estO{(s.west)}{1}{stato}
  \foreach \x in {a, b, e}{ \draw[gen] (\x.north) |- ($(s.south)+(0,-0.4)$) -- (s.south); }
\end{tikzpicture}

\medskip
\begin{tikzpicture}
  \node[iniziale] (i) at (-1.4,0) {};
  \node[stato, minimum width=1.4cm, minimum height=0.6cm] (at) at (0,0) {Attesa};
  \node[stato, minimum width=1.6cm, minimum height=0.6cm] (sc) at (3.6,0) {AttesaScelta};
  \node[stato, minimum width=1.6cm, minimum height=0.6cm] (er) at (3.6,-1.9) {Erogazione};
  \draw[tr] (i) -- (at);
  \draw[tr] ($(at.east)+(0,0.12)$) -- node[above, font=\tiny]{moneta [credito$\geq$2]} ($(sc.west)+(0,0.12)$);
  \draw[tr] ($(sc.south)+(0.3,0)$) -- node[right, font=\tiny]{scelta} ($(er.north)+(0.3,0)$);
  \draw[tr] ($(er.north)+(-0.3,0)$) -- node[left, font=\tiny, align=right]{prelievo\\{}[credito$\geq$2]} ($(sc.south)+(-0.3,0)$);
  \draw[tr] (er.west) -| node[below, font=\tiny, pos=0.3]{prelievo [credito$<$2]} (at.south);
\end{tikzpicture}

\medskip
\begin{lstlisting}[style=shellcs, name=Terminale]
$ javac Cliente.java
$ java Cliente
Credito insufficiente
Credito: 1
Credito: 2
Erogazione...
Attendi l'erogazione
Ecco il prodotto, credito: 0
\end{lstlisting}
\regola{La stessa chiamata (\code{selezionaProdotto()}) fa cose diverse a seconda dello stato: in Attesa rifiuta, in
AttesaScelta avvia l'erogazione. Una \textbf{transizione} non crea oggetti: \code{setState} sposta il riferimento
fra tre oggetti stato creati una volta sola, senza campi (il credito sta nel distributore). Il codice del corso ha
più eventi (restituzione del credito, time out, spegnimento).}
\end{programma}
